\pdfbookmark{Общая характеристика работы}{characteristic}             % Закладка pdf
\section*{Общая характеристика работы}

\newcommand{\actuality}{\pdfbookmark[1]{Актуальность}{actuality}}%\underline{\textbf{\actualityTXT}}}
\newcommand{\progress}{\pdfbookmark[1]{Разработанность темы}{progress}\underline{\textbf{\progressTXT}}}
\newcommand{\aim}{\pdfbookmark[1]{Цели}{aim}\underline{{\textbf\aimTXT}}}
\newcommand{\tasks}{\pdfbookmark[1]{Задачи}{tasks}\underline{\textbf{\tasksTXT}}}
\newcommand{\aimtasks}{\pdfbookmark[1]{Цели и задачи}{aimtasks}\aimtasksTXT}
\newcommand{\novelty}{\pdfbookmark[1]{Научная новизна}{novelty}\underline{\textbf{\noveltyTXT}}}
\newcommand{\influence}{\pdfbookmark[1]{Практическая значимость}{influence}\underline{\textbf{\influenceTXT}}}
\newcommand{\methods}{\pdfbookmark[1]{Методология и методы исследования}{methods}\underline{\textbf{\methodsTXT}}}
\newcommand{\defpositions}{\pdfbookmark[1]{Положения, выносимые на защиту}{defpositions}\underline{\textbf{\defpositionsTXT}}}
\newcommand{\reliability}{\pdfbookmark[1]{Достоверность}{reliability}\underline{\textbf{\reliabilityTXT}}}
\newcommand{\probation}{\pdfbookmark[1]{Апробация}{probation}\underline{\textbf{\probationTXT}}}
\newcommand{\contribution}{\pdfbookmark[1]{Личный вклад}{contribution}\underline{\textbf{\contributionTXT}}}
\newcommand{\publications}{\pdfbookmark[1]{Публикации}{publications}\underline{\textbf{\publicationsTXT}}}

\input{common/characteristic} % Характеристика работы по структуре во введении и в автореферате не отличается

\underline{\textbf{Объем и структура работы.}} Диссертация состоит из введения, трех глав,
заключения и одного приложения. Полный объем работы составляет 96 страниц, включая 11 рисунков и
23 таблицы. Список литературы содержит 100 наименований.

\pdfbookmark{Содержание работы}{description}                          % Закладка pdf
\section*{Содержание работы}

Во \underline{\textbf{введении}} обоснована актуальность исследования спектроскопических параметров
изотопологов молекул различной симметрии, сформулированы цель и задачи работы, определены объект и
предмет исследования, изложены научная новизна, теоретическая и практическая значимость полученных
результатов, приведены положения, выносимые на защиту, и сведения о достоверности результатов и
личном вкладе автора. Показано, что объединение экспериментальной спектроскопии высокого разрешения,
эффективных гамильтонианов и теории изотопозамещения позволяет решать обратную спектроскопическую
задачу как для сферических, так и для асимметричных волчков.

\underline{\textbf{Первая глава}} посвящена методам, составляющим теоретическую основу исследования.
Последовательно рассмотрены полный и колебательно-вращательный гамильтонианы молекулы, приближение
Борна--Оппенгеймера, теория изотопозамещения и аппарат неприводимых тензорных операторов. Исходной
точкой является стационарное уравнение Шредингера
\begin{equation}
    \widehat{H}\Psi=E\Psi,
\end{equation}
в котором точность вычисляемых энергий определяется как физической полнотой оператора
$\widehat{H}$, так и выбором базисных функций. Для описания экспериментальных уровней используется
эффективный гамильтониан, параметры которого определяются из спектроскопических данных методом
наименьших квадратов. Для молекул типа асимметричного волчка применен гамильтониан Уотсона в
$A$-редукции, содержащий вращательные постоянные и параметры центробежного искажения различных
порядков.

В приближении Борна--Оппенгеймера электронное и ядерное движения разделяются вследствие большого
различия масс электронов и ядер. Электронная потенциальная поверхность в декартовом представлении не
зависит от изотопного состава, тогда как замена ядер изменяет кинетическую энергию, моменты инерции,
нормальные координаты и частоты колебаний. Это обстоятельство лежит в основе теории
изотопозамещения: данные по хорошо изученному изотопологу используются как физические ограничения
при анализе менее изученных модификаций одной молекулы. Показано, что при симметричном замещении
сохраняется ориентация молекулярно-фиксированных осей, а при понижении симметрии, как в переходе
H$_2$S$\rightarrow$HDS, необходимо дополнительно учитывать поворот системы координат.

Аппарат неприводимых тензорных операторов используется для симметризации базисных функций,
построения матричных элементов эффективного гамильтониана и оператора дипольного момента, а также
для строгого вывода правил отбора. Для молекул с тетраэдрической симметрией T$_d$ он позволяет
разбивать гамильтониан на блоки заданной симметрии и корректно описывать тетраэдрическое расщепление
и резонансное взаимодействие вырожденных колебательных состояний. Для асимметричных волчков этот же
формализм обеспечивает согласованное описание энергетических уровней и интенсивностей разрешенных
колебательно-вращательных переходов.

\underline{\textbf{Во второй главе}} решена обратная спектроскопическая задача для молекул GeH$_4$
и SiH$_4$, относящихся к сферическим волчкам с симметрией T$_d$. В качестве исходных данных
использованы экспериментальные центры колебательных полос. Определяемыми величинами являлись
гармонические частоты, ангармонические параметры, параметры тетраэдрического расщепления и
резонансных взаимодействий. Симметризованная колебательная функция молекулы XY$_4$ представлена как
тензорное произведение функций невырожденного, дважды вырожденного и двух трижды вырожденных
осцилляторов:
\begin{equation}
\begin{aligned}
&\left|v_1;v_2l_2\Gamma_2;v_3l_3,v_4l_4,l n_l\Gamma_{34};n\Gamma\sigma\right\rangle =\\
&\qquad \left|v_1\right\rangle
\left(\left|v_2l_2\Gamma_2\right\rangle\otimes
\left|v_3l_3,v_4l_4,l n_l\Gamma_{34}\right\rangle\right)^{n\Gamma}_{\sigma}.
\end{aligned}
\label{eq:syn_basis}
\end{equation}
В таком базисе матричный элемент эффективного колебательного гамильтониана записывается в виде
\begin{equation}
H_{vv'}^{n\Gamma,n'\Gamma'}=
\widetilde{E}_{v}\delta_{vv'}\delta_{nn'}\delta_{\Gamma\Gamma'}+
W_{vv'}^{n\Gamma,n'\Gamma'},
\label{eq:syn_hamiltonian}
\end{equation}
где $\widetilde{E}_{v}$ содержит гармонические и ангармонические вклады, а
$W_{vv'}^{n\Gamma,n'\Gamma'}$ описывает тетраэдрическое расщепление при $v=v'$ и резонансные
взаимодействия при $v\ne v'$. Аналитическая форма функций и матричных элементов позволяет проводить
устойчивую подгонку параметров и затем рассчитывать центры еще не измеренных полос.

При определении параметров GeH$_4$ и SiH$_4$ учитывался локально-модовый характер растягивающих
колебаний. В предельном приближении растягивающие частоты и основные ангармонические параметры
связаны соотношениями
\begin{equation}
\omega_1=\omega+3\lambda,\qquad \omega_3=\omega-\lambda,
\end{equation}
\begin{equation}
 x_{11}\simeq\frac{1}{4}x_{13}\simeq\frac{5}{9}x_{33}\simeq
 -\frac{5}{3}G_{33}\simeq-5T_{33}\simeq\frac{1}{4}d_{1133}\simeq
 \frac{1}{16}d_{1333}.
\label{eq:syn_local}
\end{equation}
Эти зависимости применялись как физические ограничения, исключающие формально приемлемые, но плохо
интерпретируемые наборы параметров.

Для ${}^{76}$GeH$_4$ по 31 экспериментальному центру полос определен набор из 27 параметров,
воспроизводящий исходные данные со среднеквадратичным отклонением
$d_{\textrm{rms}}=0.13$ см$^{-1}$. Для шести центров полос, не включенных в подгонку, отклонение
составило 0.26 см$^{-1}$, что подтверждает предсказательную способность модели. Для изотопологов
${}^{74}$GeH$_4$, ${}^{73}$GeH$_4$, ${}^{72}$GeH$_4$ и ${}^{70}$GeH$_4$ выполнена совместная
подгонка 77 экспериментальных центров полос. С учетом регулярной масс-зависимости одноименных
параметров и близости параметров тетраэдрического расщепления получен набор из 42 варьируемых
параметров с $d_{\textrm{rms}}=0.18$ см$^{-1}$. На этой основе рассчитаны центры полос пяти
стабильных изотопологов германа для полиад с $N\leq4$~\cite{ulen1,ulen2}.

Для ${}^{28}$SiH$_4$ использованы 49 опубликованных экспериментальных центров полос и 23 центра из
базы STDS, включенные в подгонку с пониженным весом. Полученный полуэмпирический набор параметров
воспроизводит исходные центры со среднеквадратичным отклонением 0.28 см$^{-1}$. Для сравниваемого
набора, основанного на квантово-химическом силовом поле, соответствующее отклонение равно
2.21 см$^{-1}$. Таким образом, полуэмпирическое решение более чем в восемь раз точнее описывает
известные центры полос и одновременно лучше удовлетворяет локально-модовым соотношениям. Для
${}^{29}$SiH$_4$ и ${}^{30}$SiH$_4$ совместно обработаны 32 центра полос. При подгонке использовано
приближенное изотопическое условие
\begin{equation}
{}^{29}P_n\simeq\frac{1}{2}\left({}^{28}P_n+{}^{30}P_n\right),
\end{equation}
что позволило определить 18 варьируемых параметров с $d_{\textrm{rms}}=0.06$ см$^{-1}$. Итоговые
наборы параметров использованы для расчета 6375 центров колебательных полос до 9000 см$^{-1}$, то
есть по 2125 центров для каждого из трех стабильных изотопологов SiH$_4$.

\begin{table}[ht]
\centering
\caption{Основные результаты решения обратной задачи для молекул типа XY$_4$(T$_d$).}
\label{tab:syn_xy4_results}
\begingroup
\scriptsize
\setlength{\tabcolsep}{3pt}
\setlength{\extrarowheight}{2pt}
\renewcommand{\arraystretch}{1.12}
\begin{tabularx}{\textwidth}{|>{\centering\arraybackslash}p{0.20\textwidth}|>{\centering\arraybackslash}p{0.18\textwidth}|>{\centering\arraybackslash}p{0.18\textwidth}|X|}
\hline
Объект & Исходные данные & $d_{\textrm{rms}}$, см$^{-1}$ & Итог анализа \\
\hline
${}^{76}$GeH$_4$ & 31 центр полос & 0.13 & 27 параметров; расчет уровней при $N\leq4$ \\
\hline
${}^{74,73,72,70}$GeH$_4$ & 77 центров полос & 0.18 & 42 параметра; единое описание четырех изотопологов \\
\hline
${}^{28}$SiH$_4$ & 49+23 центра полос & 0.28 & 2125 рассчитанных центров до 9000 см$^{-1}$ \\
\hline
${}^{29,30}$SiH$_4$ & 32 центра полос & 0.06 & 18 параметров; по 2125 центров для каждого изотополога \\
\hline
\end{tabularx}
\endgroup
\end{table}

\underline{\textbf{Третья глава}} посвящена анализу высокоразрешенных спектров дейтерированных
изотопологов сероводорода D$_2{}^{32}$S, HD$^{32}$S и HD$^{34}$S в области фундаментальной
деформационной полосы $\nu_2$. Спектры зарегистрированы на Фурье-спектрометре Bruker IFS~125HR с
разрешением 0.002 и 0.003 см$^{-1}$ при температуре $23\pm0.8\,{}^\circ$C и полном давлении
450 Па. Использованы многоходовые кюветы с оптическими длинами пути 4 и 163 м. Для D$_2$S
анализировалась область 750--1200 см$^{-1}$, для HDS --- 690--1510 см$^{-1}$. Короткая кювета
использовалась при определении интенсивностей, поскольку она уменьшала насыщение сильных линий.
Площади изолированных линий находились из подгонки профиля Хартмана--Трана, а интенсивности
рассчитывались по закону Бугера--Ламберта--Бера:
\begin{equation}
S=\frac{1}{\lg e}\frac{k_BT}{P L}
\int\lg\left(\frac{I_0}{I}\right)_{\widetilde{\nu}}d\widetilde{\nu}.
\label{eq:syn_intensity}
\end{equation}

Молекула D$_2$S имеет симметрию C$_{2v}$, а ее полоса $\nu_2$ относится к полосам $b$-типа. Линии
относились методом комбинационных разностей. Выявлен систематический рост отклонений литературных
энергий основного состояния при больших $J$ и $K_a$, поэтому в подгонку были дополнительно включены
270 экспериментальных комбинационных разностей. Уточненный набор параметров воспроизводит их с
$d_{\textrm{rms}}=1.4\times10^{-4}$ см$^{-1}$, что в 2.9 раза лучше прежнего описания. После этого
к полосе $\nu_2$ D$_2{}^{32}$S отнесено 1742 перехода с $J_{\max}=30$ и
$K_{a,\max}=21$. По ним определено 535 энергий состояния $(010)$. Подгонка 34 параметров
гамильтониана Уотсона воспроизводит эти энергии со среднеквадратичным отклонением
$1.31\times10^{-4}$ см$^{-1}$~\cite{d2s_sydow2019}.

Вследствие водород-дейтериевого обмена исходный образец D$_2$S в ходе длительной регистрации
превращался в смесь D$_2$S, HDS и H$_2$S, поэтому парциальное давление D$_2$S нельзя было принять
равным полному давлению образца. Для его определения из параметров дипольного момента H$_2$S и
соотношений теории изотопозамещения получена оценка главного параметра полосы
${}^{010}\widetilde{\mu}_{x1}=0.008519$ D. Экспериментальные и рассчитанные интенсивности пробных
линий связывались формулой
\begin{equation}
P_{\textrm{part}}=
\frac{{}^{(\textrm{prob})}S^N_{\nu}}{{}^{(\textrm{calc})}S^N_{\nu}}
P_{\textrm{samp}}.
\label{eq:syn_pressure}
\end{equation}
Учет двух главных параметров эффективного дипольного момента дал
$P_{\textrm{part}}=2.378\pm0.036$ мбар при полном давлении 4.5 мбар. После введения этой поправки
проанализированы интенсивности 280 линий, соответствующих 324 переходам с $J_{\max}=16$ и
$K_{a,\max}=9$. Набор из восьми параметров эффективного оператора дипольного момента воспроизводит
экспериментальные интенсивности с $d_{\textrm{rms}}=2.5$~\%; итоговая неопределенность рассчитанных
значений оценена приблизительно в 3.5~\%.

Для HDS понижение симметрии от C$_{2v}$ к C$_s$ приводит к гибридному характеру полос и требует
учета поворота молекулярно-фиксированной системы координат при выводе изотопических соотношений.
Для HD$^{32}$S к полосе $\nu_2$ отнесено 2684 перехода с $J_{\max}=25$ и
$K_{a,\max}=17$, по которым получено 415 энергий состояния $(010)$. Набор из 33 параметров
гамильтониана Уотсона описывает их с $d_{\textrm{rms}}=1.1\times10^{-4}$ см$^{-1}$. Для менее
распространенного изотополога HD$^{34}$S отнесено 1212 переходов с $J_{\max}=21$ и
$K_{a,\max}=14$, определено 263 энергии состояния $(010)$ и достигнуто отклонение
$1.6\times10^{-4}$ см$^{-1}$~\cite{hds_sydow2019}.

Из изотопических соотношений получены оценки главных параметров эффективного дипольного момента HDS
${}^{010}\widetilde{\mu}_{x1}=0.01886$ D и
${}^{010}\widetilde{\mu}_{z1}=-0.00855$ D. Сопоставление рассчитанных и экспериментальных
интенсивностей десяти пробных линий дало парциальное давление HD$^{32}$S
$P_{\textrm{part}}=1.77\pm0.10$ мбар, соответствующее доле около 39.4~\% в образце. После этого
определены интенсивности 418 линий с $J_{\max}=21$ и $K_{a,\max}=6$. Полученный набор параметров
эффективного дипольного момента воспроизводит их с $d_{\textrm{rms}}=4.7$~\%. Отношение главных
параметров, найденное из подгонки, равно 2.183 и согласуется с теоретическим значением 2.20. Для
наиболее сильных линий HD$^{34}$S соответствующее отклонение составляет около 4.0~\%.

\begin{table}[ht]
\centering
\caption{Основные результаты анализа полосы $\nu_2$ дейтерированных изотопологов сероводорода.}
\label{tab:syn_h2s_results}
\begingroup
\scriptsize
\setlength{\tabcolsep}{2.5pt}
\renewcommand{\arraystretch}{1.12}
\begin{tabularx}{\textwidth}{|>{\centering\arraybackslash}p{0.15\textwidth}|>{\centering\arraybackslash}p{0.16\textwidth}|>{\centering\arraybackslash}p{0.16\textwidth}|>{\centering\arraybackslash}p{0.18\textwidth}|X|}
\hline
Молекула & Переходы & Энергии $(010)$ & $d_{\textrm{rms}}$ энергий, см$^{-1}$ & Анализ интенсивностей \\
\hline
D$_2{}^{32}$S & 1742 & 535 & $1.31\times10^{-4}$ & 280 линий; 2.5~\% \\
\hline
HD$^{32}$S & 2684 & 415 & $1.1\times10^{-4}$ & 418 линий; 4.7~\% \\
\hline
HD$^{34}$S & 1212 & 263 & $1.6\times10^{-4}$ & сильные линии; около 4.0~\% \\
\hline
\end{tabularx}
\endgroup
\end{table}

В результате исследования сформированы согласованные списки положений и рассчитанных интенсивностей
1742 переходов D$_2{}^{32}$S, 2684 переходов HD$^{32}$S и 1212 переходов HD$^{34}$S. Полученные
данные описывают энергетическую структуру и абсолютные интенсивности фундаментальной полосы $\nu_2$
с точностью, близкой к экспериментальной, и могут использоваться при моделировании спектров высокого
разрешения, уточнении потенциальной и дипольной поверхностей, а также при подготовке данных для
прикладных задач атмосферной оптики, астрофизики и планетологии.

\FloatBarrier
\pdfbookmark{Заключение}{conclusion}                                  % Закладка pdf
В \underline{\textbf{заключении}} приведены основные результаты работы, которые заключаются в следующем:
\input{common/concl}

\pdfbookmark{Литература}{bibliography}                                % Закладка pdf
\ifdefmacro{\microtypesetup}{\microtypesetup{protrusion=false}}{} % не рекомендуется применять пакет микротипографики к автоматически генерируемому списку литературы
\urlstyle{rm}                               % ссылки URL обычным шрифтом
\ifnumequal{\value{bibliosel}}{0}{% Встроенная реализация с загрузкой файла через движок bibtex8
    \renewcommand{\bibname}{\large \bibtitleauthor}
    \nocite{*}
    \insertbiblioauthor           % Подключаем Bib-базы
    %\insertbiblioexternal   % !!! bibtex не умеет работать с несколькими библиографиями !!!
}{% Реализация пакетом biblatex через движок biber
    % Цитирования.
    %  * Порядок перечисления определяет порядок в библиографии (только внутри подраздела, если `\insertbiblioauthorgrouped`).
    %  * Если не соблюдать порядок "как для \printbibliography", нумерация в `\insertbiblioauthor` будет кривой.
    %  * Если цитировать каждый источник отдельной командой --- найти некоторые ошибки будет проще.
    %
    %% authorvak
    \nocite{vakbib1}%
    \nocite{vakbib2}%
    \nocite{vakbib3}%
    \nocite{vakbib4}%
    \nocite{vakbib5}%
    \nocite{vakbib6}%
    \nocite{vakbib7}%
    \nocite{vakbib8}%
    \nocite{vakbib9}%
    \nocite{vakbib10}%
    \nocite{vakbib11}%
    \nocite{vakbib12}%
    %
    %% authorwos
    \nocite{wosbib1}%
    %
    %% authorscopus
    \nocite{scbib1}%
    %
    %% authorpatent
    \nocite{patbib1}%
    %
    %% authorprogram
    \nocite{progbib1}%
    %
    %% authorconf
    \nocite{confbib1}%
    \nocite{confbib2}%
    %
    %% authorother
    \nocite{bib1}%
    \nocite{bib2}%

    \ifnumgreater{\value{usefootcite}}{0}{
        \begin{refcontext}[labelprefix={}]
            \ifnum \value{bibgrouped}>0
                \insertbiblioauthorgrouped    % Вывод всех работ автора, сгруппированных по источникам
            \else
                \insertbiblioauthor      % Вывод всех работ автора
            \fi
        \end{refcontext}
    }{
        \ifnum \totvalue{citeexternal}>0
            \begin{refcontext}[labelprefix=A]
                \ifnum \value{bibgrouped}>0
                    \insertbiblioauthorgrouped    % Вывод всех работ автора, сгруппированных по источникам
                \else
                    \insertbiblioauthor      % Вывод всех работ автора
                \fi
            \end{refcontext}
        \else
            \ifnum \value{bibgrouped}>0
                \insertbiblioauthorgrouped    % Вывод всех работ автора, сгруппированных по источникам
            \else
                \insertbiblioauthor      % Вывод всех работ автора
            \fi
        \fi
        %  \insertbiblioauthorimportant  % Вывод наиболее значимых работ автора (определяется в файле characteristic во второй section)
        \begin{refcontext}[labelprefix={}]
            \insertbiblioexternal            % Вывод списка литературы, на которую ссылались в тексте автореферата
        \end{refcontext}
        % Невидимый библиографический список для подсчёта количества внешних публикаций
        % Используется, чтобы убрать приставку "А" у работ автора, если в автореферате нет
        % цитирований внешних источников.
        \printbibliography[heading=nobibheading, section=0, env=countexternal, keyword=biblioexternal, resetnumbers=true]%
    }
}
\ifdefmacro{\microtypesetup}{\microtypesetup{protrusion=true}}{}
\urlstyle{tt}                               % возвращаем установки шрифта ссылок URL
